\documentclass[12pt]{article}

\usepackage[a4paper, margin=2.5cm]{geometry}
\usepackage{amsmath}
\usepackage{amssymb}
\usepackage{amsthm}
\usepackage[colorlinks=true, linkcolor=blue, citecolor=blue, urlcolor=blue]{hyperref}
\usepackage{booktabs}
\usepackage{natbib}

\newtheorem{theorem}{Theorem}
\newtheorem{proposition}[theorem]{Proposition}
\newtheorem{corollary}[theorem]{Corollary}
\newtheorem{definition}[theorem]{Definition}
\newtheorem{remark}[theorem]{Remark}
\newtheorem{example}[theorem]{Example}
\newtheorem{lemma}[theorem]{Lemma}

\newcommand{\Ecal}{\mathcal{E}}
\newcommand{\Gcal}{\mathcal{G}}
\newcommand{\Fcal}{\mathcal{F}}
\newcommand{\Ocal}{\mathcal{O}}
\newcommand{\Ccal}{\mathcal{C}}
\newcommand{\Pst}{P_{\mathrm{st}}}

\title{Measurement as a Conditional Barrier Template\\for Sector Separation in Closure Dynamics}
\author{%
  Lee Smart\\[4pt]
  \textit{Institute of Vibrational Field Dynamics}\\[2pt]
  \texttt{contact@vibrationalfielddynamics.org}\\[2pt]
  \texttt{@vfd\_org}%
}
\date{April 2026}

\begin{document}

\maketitle

\noindent\textbf{Status:} Preprint (not peer-reviewed)

\begin{abstract}
Paper~XXX proposed that outcome probabilities arise as geometric weights over observer-admissible sectors, with the Born rule recovered from the stationary sector measure of the closure dynamics under five assumptions: equilibrium, well-separated sectors, sector-supported decomposition, a sampling postulate, and state-completeness. The present paper addresses, at most, one possible route to the metastability part of the well-separated-sectors assumption, by isolating a conditional sufficient-condition template: \textit{under what further hypotheses might inter-outcome barriers be large in a measurement context?} The negligible-overlap part of Paper~XXX's assumption~2, and the sector-supported decomposition of assumption~3, are not established here.

The proposal is that measurement interactions may dynamically reshape the closure landscape, creating or correlating apparatus-correlated basins corresponding to distinct outcomes (depending on whether the pre-measurement landscape is mono- or multi-stable). Concretely, system--apparatus coupling introduces an interaction term $\Fcal_{\mathrm{int}}$ into the total closure functional, and \emph{if} the post-interaction landscape is quasi-static and each relevant escape can be reduced to a Paper~XIV-type two-basin saddle problem, the imported small-noise Kramers asymptotic would predict exponentially suppressed transitions. Whether such deformation actually occurs for physical measurement couplings, and whether the resulting landscape yields a genuine sector-supported decomposition, are separate questions not settled here.

A toy model illustrates the mechanism class. Its algebraic details are worked out in Section~\ref{sec:toy}: within the parameter regime where the bistable structure survives, the barrier depends monotonically on coupling, but the sign of that dependence is model-specific and does not, by itself, establish a generic timescale hierarchy. A heuristic timescale ansatz suggests that sector formation \emph{can} occur faster than measurement readout in a suitably chosen regime; no theorem to that effect is proved.

The paper does not claim a complete solution to the measurement problem, does not discharge Paper~XXX's Route~C assumptions (four of the five remain external), and does not replace a decoherence account. It identifies a sufficient-condition template and points to the quantitative work needed to turn it into a derivation.
\end{abstract}

%==================================================================
\section{Introduction}\label{sec:intro}
%==================================================================

Paper~XXX~\citep{SmartXXX} showed that, in the equilibrium closure regime, the Born rule can be recovered from the stationary sector measure of the closure dynamics. The argument (Paper~XXX's Proposition on Conditional Born Weighting from Stationary Sector Measure) depends on five assumptions:

\begin{enumerate}
    \item The system is at the Paper~XIV equilibrium ($\rho = \Pst$).
    \item The outcome sectors are well-separated basins with negligible overlap.
    \item The Nelson wavefunction admits a decomposition into sector-supported components.
    \item A sampling postulate: measurement frequencies equal integrated stationary sector masses.
    \item State-completeness: the Nelson density is supported in the observer-admissible set.
\end{enumerate}

Paper~XIV establishes a stationary distribution but does \emph{not} prove generic convergence of arbitrary initial data to equilibrium; assumption~1 therefore remains conditional. Assumption~3 requires not only disjoint basins but an actual spectral construction of sector-localised modes. Assumptions~4 and~5 are measurement-theoretic/observer-side postulates that the present paper does not address. The present paper addresses only one possible route to the \emph{metastability component} of assumption~2: \textit{under what conditions could inter-outcome barriers be large in a measurement setup?} The negligible-overlap part of assumption~2, and the sector-supported decomposition of assumption~3, are not established here.

\subsection*{The gap}

In standard quantum mechanics, the transition from superposition to definite outcomes is typically attributed to decoherence: interaction with the environment suppresses off-diagonal elements of the density matrix. The VFD framework does not use density matrices or environmental tracing. The question is whether the closure dynamics provides an intrinsic mechanism for sector separation.

\subsection*{Proposal}

This paper conditionally repackages sector separation into four explicit additional hypotheses on the system--apparatus landscape (Remark~\ref{rem:basins} below), under which measurement-type coupling would yield a pairwise barrier lower bound. The framing is geometric: the interaction term is \emph{intended} to reshape the landscape into multiple basin regions separated by barriers, but the paper does not derive the required topological/gap structure from the measurement coupling alone. A toy model is examined whose bistable regime \emph{fails} to exhibit coupling-driven barrier amplification.

\subsection*{Scope}

The paper develops the formal framework, works a toy model, provides a heuristic timescale ansatz, and connects the mechanism class to Paper~XXX. It does not claim a complete solution to the measurement problem, does not derive decoherence from the closure framework, does not specify the interaction functional from first principles, and does not prove that physical measurement interactions generate the sector structure Paper~XXX requires. It identifies one candidate mechanism class that \emph{could} produce the metastability component of well-separated sectors in a measurement context, and flags the quantitative work needed to discharge Paper~XXX's assumption~2 rigorously (including the overlap/negligible-interference component, which this paper does not settle).

%==================================================================
\section{Programme Position}\label{sec:position}
%==================================================================

The logical chain:

\begin{enumerate}
    \item \textbf{Structure}: closure geometry on the 600-cell (Papers~I--V, XXII).
    \item \textbf{Dynamics}: stochastic closure process with stationary distribution (Paper~XIV).
    \item \textbf{Quantum recovery}: Schr\"odinger equation as equilibrium tangent limit (Papers~XVIII--XXI).
    \item \textbf{Observer}: constraint substructure determining admissible configurations (Paper~XXIX).
    \item \textbf{Measurement}: dynamical sector separation via landscape deformation (\textit{this paper}).
    \item \textbf{Probability}: conditional stationary sector measure yields Born weighting under Paper~XXX's five Route~C hypotheses (Paper~XXX).
\end{enumerate}

Steps 5 and 6 were reversed in the order of publication: Paper~XXX appeared first because the probability construction is logically downstream, but it depends on a sector-separation assumption whose mechanism is explored here. This paper narrows the gap between steps 4 and 6 by isolating a candidate mechanism class and the missing hypotheses it rests on; it does not fill that gap.

%==================================================================
\section{Formal Setup}\label{sec:setup}
%==================================================================

\subsection{Configuration space}

Let $\Phi_S \in \mathcal{H}_S$ denote the system configuration and $\Phi_A \in \mathcal{H}_A$ the apparatus configuration, where $\mathcal{H}_S$ and $\mathcal{H}_A$ are the respective configuration spaces. The total configuration is $\Phi = (\Phi_S, \Phi_A) \in \mathcal{H}_S \times \mathcal{H}_A$.

\subsection{Closure functional}

\begin{definition}[Total closure functional]\label{def:total}
The total closure functional for a measurement context is:
\begin{equation}\label{eq:total}
    \Fcal_{\mathrm{total}}[\Phi_S, \Phi_A] = \Fcal_S[\Phi_S] + \Fcal_A[\Phi_A] + \Fcal_{\mathrm{int}}[\Phi_S, \Phi_A]
\end{equation}
where $\Fcal_S$ is the system closure functional, $\Fcal_A$ is the apparatus closure functional, and $\Fcal_{\mathrm{int}}$ is the interaction term encoding the measurement coupling.
\end{definition}

\begin{remark}
Before measurement, $\Fcal_{\mathrm{int}} = 0$ and the system and apparatus evolve independently. Measurement begins when the interaction is switched on ($\Fcal_{\mathrm{int}} \neq 0$). The interaction term does not destroy the closure structure; it deforms the landscape.
\end{remark}

\subsection{Stochastic dynamics}

The total configuration evolves under the stochastic closure dynamics (Paper~XIV):
\begin{equation}\label{eq:dynamics}
    d\Phi = -\nabla \Fcal_{\mathrm{total}}\, dt + \sigma\, dW
\end{equation}
where $\sigma > 0$ is the noise strength and $dW$ is a Wiener process on the relevant finite-dimensional joint configuration model. In a finite-dimensional joint configuration model satisfying the Paper~XIV smoothness/confinement hypotheses, the stationary density would be
\begin{equation}\label{eq:stationary}
    \Pst(\Phi) \propto e^{-2\Fcal_{\mathrm{total}}(\Phi)/\sigma^2}.
\end{equation}
We use this form throughout. Uniqueness of this stationary distribution is conditional on the irreducibility/regularity hypotheses of Paper~XIV; generic convergence from arbitrary initial data is not established there. An infinite-dimensional SPDE extension would require a separate stationary-measure construction and is not carried out here.

%==================================================================
\section{Pre-Measurement Landscape}\label{sec:pre}
%==================================================================

Before the interaction is switched on ($\Fcal_{\mathrm{int}} = 0$), the total functional is separable:
\begin{equation}
    \Fcal_{\mathrm{total}} = \Fcal_S[\Phi_S] + \Fcal_A[\Phi_A]
\end{equation}

The pre-measurement landscape depends on the specific $\Fcal_S$ and $\Fcal_A$ chosen. For a bistable system (e.g.\ the double well of Section~\ref{sec:toy}), $\Fcal_S$ already has two minima, and the uncoupled total landscape has two minima at $(\phi_S^{(i)}, \Phi_A^\star)$ where $\Phi_A^\star$ is the apparatus ready state; crucially, in the uncoupled case the apparatus configuration does \emph{not} distinguish outcomes. The interest of the post-measurement landscape is therefore \emph{apparatus correlation} with outcome, not necessarily the appearance of new minima. For a non-bistable system (a single broad minimum of $\Fcal_S$), the uncoupled landscape has a single basin and coupling is required to produce multiple minima; that is the traditional ``measurement induces branches'' picture. The present paper is not committed to either sub-case; the mechanism it studies is the coupling-induced deformation of the landscape, whatever its pre-measurement structure.

\begin{quote}
Pre-measurement: apparatus configuration is uncorrelated with outcome, regardless of whether $\Fcal_S$ itself is mono- or multi-stable.
\end{quote}

%==================================================================
\section{Measurement Interaction and Landscape Deformation}\label{sec:deformation}
%==================================================================

When the measurement interaction is switched on, $\Fcal_{\mathrm{int}}$ couples the system and apparatus configurations. The interaction term introduces correlations: the closure cost of the total configuration depends on the joint state of system and apparatus.

\subsection{Correlation structure}

A measurement interaction is characterised by the property that it energetically favours configurations in which the apparatus state is correlated with the system state. Formally:

\begin{definition}[Measurement-type interaction]\label{def:measurement}
An interaction $\Fcal_{\mathrm{int}}[\Phi_S, \Phi_A]$ is of \textit{measurement type} if:
\begin{enumerate}
    \item It has $N$ low-cost channels: for each designated system configuration $\Phi_S^{(i)}$ (a chosen local minimum of $\Fcal_S$, or the closure-space counterpart of a system eigenstate where applicable), there exists a corresponding apparatus configuration $\Phi_A^{(i)}$ such that $\Fcal_{\mathrm{int}}[\Phi_S^{(i)}, \Phi_A^{(i)}]$ is small.
    \item Mismatched configurations have high interaction cost: $\Fcal_{\mathrm{int}}[\Phi_S^{(i)}, \Phi_A^{(j)}] \gg \Fcal_{\mathrm{int}}[\Phi_S^{(i)}, \Phi_A^{(i)}]$ for $i \neq j$.
\end{enumerate}
\end{definition}

\begin{remark}
This is the closure-dynamical analogue of the standard requirement that a measurement apparatus correlates its pointer states with the system's eigenstates. The VFD version is expressed in terms of the closure landscape rather than Hilbert-space operators.
\end{remark}

\subsection{Landscape deformation}

The effect of a measurement-type interaction on the total closure landscape is:

\begin{remark}[Heuristic template for sector-forming measurement interactions]\label{rem:basins}
We record a sufficient-condition \emph{template}, not a proved theorem. Crucially, hypotheses (ii) (topological separation) and (iii) (positive barrier gap) below already encode most of the separation content the template appears to deliver: the paper does \emph{not} derive metastability from measurement coupling, but assumes a separation topology plus a quantitative gap and then records the immediate barrier consequence. The contribution is an explicit accounting of what is assumed versus what follows, not a mechanism. Suppose $\Fcal_{\mathrm{int}}$ is of measurement type (Definition~\ref{def:measurement}) and the correlated channels $\Phi^{(i)*} := (\Phi_S^{(i)}, \Phi_A^{(i)})$ are nondegenerate local minima of $\Fcal_{\mathrm{total}}$ under standard regularity assumptions. For each ordered pair $(i,j)$ define the directional barrier
\[
    \Delta\Fcal_{i\to j} \;:=\; \inf_{\gamma:\,i \to j}\,\sup_{\Phi \in \gamma}\, \Fcal_{\mathrm{total}}(\Phi) \;-\; \Fcal_{\mathrm{total}}(\Phi^{(i)*}),
\]
where the infimum runs over continuous paths from $\Phi^{(i)*}$ to $\Phi^{(j)*}$. The template invokes four additional hypotheses beyond Definition~\ref{def:measurement}:

\emph{Hypothesis (i):} these $N$ correlated pairs exhaust the local minima of $\Fcal_{\mathrm{total}}$.

\emph{Hypothesis (iv) [Neighbourhood continuity/gap, beyond Definition~\ref{def:measurement}]:} continuity of $\Fcal_{\mathrm{int}}$ plus a quantitative gap permits disjoint open matched neighbourhoods $U_k \ni \Phi_S^{(k)}$, $V_k \ni \Phi_A^{(k)}$ on which $\Fcal_{\mathrm{int}}$ remains small on $U_k \times V_k$ and large on the cross-mismatched products $U_i \times V_j$, $U_j \times V_i$. Definition~\ref{def:measurement} controls only the discrete designated points; neighbourhood-level control requires this additional continuity/gap assumption. Define the pairwise cross-mismatch region
\[
    \mathrm{MM}_{ij} \;:=\; (U_i \times V_j) \,\cup\, (U_j \times V_i).
\]
\emph{Hypothesis (ii) is a strong topological separation assumption not derived here:} every continuous path $\gamma: \Phi^{(i)*} \to \Phi^{(j)*}$ must enter $\mathrm{MM}_{ij}$, i.e.\ there exists $t \in (0,1)$ with $\gamma(t) \in \mathrm{MM}_{ij}$. By construction $\mathrm{MM}_{ij}$ is disjoint from both matched minima $\Phi^{(i)*}, \Phi^{(j)*}$, so the hypothesis is non-vacuous; nothing in ``measurement-type coupling'' (Definition~\ref{def:measurement}) forces it, so (ii) largely re-imposes sector separation rather than explaining it; and (iii) the pairwise positive gap
\begin{equation}\label{eq:barrier-hyp}
    \inf_{\Phi \in \mathrm{MM}_{ij}} \Fcal_{\mathrm{total}}(\Phi) \;-\; \Fcal_{\mathrm{total}}(\Phi^{(i)*}) \;\geq\; \beta_{ij} \;>\; 0
\end{equation}
holds for each $(i,j)$, where $\beta_{ij}$ is a parameter of the model. Under (i), (ii), (iii), and (iv), the directional barrier inherits this bound directly:
\begin{equation}\label{eq:barrier}
    \Delta\Fcal_{i\to j} \;\geq\; \beta_{ij}.
\end{equation}
Whether $\beta_{ij}$ can in turn be related to raw mismatch interaction values --- e.g.\ whether $\beta_{ij} \geq \min(\Fcal_{\mathrm{int}}[\Phi_S^{(i)}, \Phi_A^{(j)}], \Fcal_{\mathrm{int}}[\Phi_S^{(j)}, \Phi_A^{(i)}]) - c_{S,A}(i,j)$ for some comparison term $c_{S,A}(i,j)$ controlling the $\Fcal_S + \Fcal_A$ contribution along connecting paths --- is a separate case-by-case modelling question. Definition~\ref{def:measurement} fixes mismatch cost only at the two discrete designated points, not over the whole mismatch set, so the relation $\beta_{ij} \mapsto $ (raw mismatch costs) is not automatic. The notation $\Delta\Fcal_{ij}$ in the rest of this paper refers to the directional barrier $\Delta\Fcal_{i\to j}$ when a direction matters, in particular in the Kramers expression below. Existence of exactly $N$ basins, existence and location of saddles between them, the uniform lower bound, the neighbourhood-level continuity/gap, and the consequent barrier estimate are not established here; (i)--(iv) are additional hypotheses that would need to be verified case by case. A rigorous statement would also require specifying the topology of the configuration space and verifying transversality/Morse-theoretic conditions on the critical points.
\end{remark}

\begin{quote}
Template (not theorem): under the four additional hypotheses (i)--(iv) above, measurement-type coupling provides a conditional repackaging of sector separation with the barrier bound~\eqref{eq:barrier}. Hypothesis (ii) is itself a strong topological separation assumption; hypotheses (iii) and (iv) are quantitative gap conditions not derived from Definition~\ref{def:measurement}. Whether a given physical coupling satisfies all four is a separate modelling question not settled here.
\end{quote}

%==================================================================
\section{Transition Suppression}\label{sec:suppression}
%==================================================================

\subsection{Kramers-type barrier crossing}

Paper~XIV~\citep{SmartXIV} states the standard small-noise Kramers asymptotic for a fixed two-basin saddle setting (citing the classical Freidlin--Wentzell framework), and we import that statement here. Assuming (a) the post-interaction landscape has frozen (the quasi-static hypothesis: landscape deformation is slow on the timescale of basin relaxation), (b) the analogous saddle/nondegeneracy hypotheses hold for each pairwise escape problem $i \to j$ in the multi-basin landscape, and (c) a justified two-basin reduction isolates a single dominant index-1 saddle controlling each pairwise escape, one would then be able to import that estimate to the directional barriers $\Delta\Fcal_{i\to j}$:
\begin{equation}\label{eq:kramers}
    \Gamma_{i \to j} \sim e^{-2\Delta\Fcal_{i \to j}/\sigma^2}.
\end{equation}
When $\Delta\Fcal_{i \to j} \gg \sigma^2$, the transition rate is exponentially suppressed. Without the quasi-static hypothesis, a time-dependent Eyring--Kramers statement would be required, which is not supplied here.

\begin{definition}[Sector separation condition]\label{def:separation}
The outcome sectors are \textit{well-separated} if the directional barriers satisfy
\begin{equation}\label{eq:separation}
    \Delta\Fcal_{i\to j} \gg \sigma^2 \qquad \text{for every ordered pair } i \neq j.
\end{equation}
\end{definition}

\begin{remark}
Within the conditional template of Remark~\ref{rem:basins}, this separation condition is what one would expect to emerge if the four hypotheses (i), (ii), (iii), and (iv) all hold; it is not derived from Definition~\ref{def:measurement} alone. In particular, ``mismatch costs are large'' at the designated discrete points is weaker than hypothesis (iii) (uniform positive gap over the mismatch set) and weaker than hypothesis (iv) (neighbourhood-level continuity/gap). Whether a given physical coupling satisfies (iii) and (iv) is a separate modelling question.
\end{remark}

\subsection{Cross-term suppression (heuristic)}

If well-separated basins \emph{and} a corresponding sector-localised basis $\{\phi_i\}$ are both available --- neither is proved in this paper --- one expects the overlap between $\phi_k$ and $\phi_l$ to be exponentially small in some tunnelling exponent, by standard semiclassical/WKB arguments. The schematic estimate
\begin{equation}\label{eq:overlap}
    |\langle \phi_k | \phi_l \rangle| \;\lesssim\; e^{-c\,\Delta\Fcal_{kl}/\sigma^2}
\end{equation}
(for some constant $c > 0$ depending on the spectral construction used) captures this expectation; we do \emph{not} derive it here. A rigorous bound would require (a) a proof that sector-localised modes $\phi_i$ exist with the support properties used by Paper~XXX's Proposition on Conditional Born Weighting, and (b) a WKB/Agmon-type tunnelling estimate. Both are deferred. Consequently, this section does \emph{not} establish Paper~XXX's assumption~3 (sector-supported decomposition).

\begin{quote}
Heuristic: interference between sectors is expected to be suppressed by the same mechanism that suppresses transitions (the barrier height), but quantifying this requires a spectral/semiclassical analysis not carried out here.
\end{quote}

%==================================================================
\section{Outcome Stability}\label{sec:stability}
%==================================================================

Once the system has settled into basin $i$, one expects it to remain there with high probability under suitable conditions. Two distinct ingredients enter, and they belong to different categories:

\begin{enumerate}
    \item \textbf{Dynamical metastability (conditional).} If all three Paper~XIV import hypotheses of Section~\ref{sec:suppression} hold (quasi-static post-interaction landscape, pairwise saddle/nondegeneracy conditions, and a justified two-basin reduction) and $\Delta\Fcal_{i\to j} \gg \sigma^2$, the Kramers rate $\Gamma_{i \to j} \sim e^{-2\Delta\Fcal_{i\to j}/\sigma^2}$ is exponentially small, so spontaneous transitions out of the basin are then extremely rare. This is the only source of dynamical metastability on offer.

    \item \textbf{Observer-side admissibility filter (separate category).} The observer $(\Ocal, \Ccal_\Ocal)$ of Paper~XXIX~\citep{SmartXXIX} defines an admissibility set $\Phi_{\mathrm{adm}}(\Ocal) = \{\Phi : \Ccal_\Ocal(\Phi) \le \epsilon\}$. If transition regions happen to fall outside $\Phi_{\mathrm{adm}}(\Ocal)$ (a case-by-case modelling question, not a theorem of Paper~XXIX), observer admissibility may separately filter which transition-region configurations count as registered outcomes. This does \emph{not} add a barrier term to $\Fcal_{\mathrm{total}}$ or by itself increase residence times; it is a category distinct from dynamical metastability.
\end{enumerate}

\begin{quote}
In the intended mechanism class, classical definiteness of measurement outcomes would correspond to stable basin occupancy in the deformed closure landscape.
\end{quote}

%==================================================================
\section{Timescale Argument}\label{sec:timescale}
%==================================================================

For the sector-separation mechanism to produce definite outcomes, the basins must form before the measurement result is read.

\begin{definition}[Separation timescale]\label{def:timescale}
The \textit{separation timescale} $\tau_{\mathrm{sep}}$ is the time required for the closure landscape to deform from an apparatus-uncorrelated configuration to one with apparatus-correlated metastable sectors after the interaction $\Fcal_{\mathrm{int}}$ is switched on. (Whether the pre-measurement landscape is mono- or multi-stable is model-dependent, as noted in Section~\ref{sec:pre}.) As a purely schematic scaling relation, assuming the dynamics has been nondimensionalised (mobility/friction coefficient normalised to unity), one may estimate:
\begin{equation}\label{eq:timescale}
    \tau_{\mathrm{sep}} \sim \frac{1}{\|\nabla^2 \Fcal_{\mathrm{int}}\|},
\end{equation}
the inverse of the curvature scale set by the interaction. In dimensional form a mobility $\mu$ and $k_B T$-like noise factor would reinstate the missing units; we do not track them here.
\end{definition}

\begin{remark}[Timescale hierarchy]
For the mechanism to work, the following hierarchy must hold:
\begin{equation}\label{eq:hierarchy}
    \tau_{\mathrm{sep}} \ll \tau_{\mathrm{relax}} \ll \tau_{\mathrm{obs}}
\end{equation}
where $\tau_{\mathrm{relax}}$ is the time for the stochastic dynamics to settle into a basin (set by the local relaxation rate of $\Fcal_{\mathrm{total}}$) and $\tau_{\mathrm{obs}}$ is the timescale on which the observer reads the outcome. The first inequality says basins form before the system relaxes into one; the second says the system has settled before the result is recorded. This hierarchy is treated as a working hypothesis for macroscopic measurement apparatus; whether it holds universally is an open question.
\end{remark}

%==================================================================
\section{Heuristic Status of the Basin-Formation Argument}\label{sec:heuristic}
%==================================================================

The present paper identifies a \textit{mechanism class} for sector formation, not a theorem. The argument identifies additional hypotheses (Remark~\ref{rem:basins}, points (i)--(iv)) under which measurement-type interactions (Definition~\ref{def:measurement}) \emph{would} produce separated basins with the barrier lower bound of~\eqref{eq:barrier}; the present paper does not establish that those hypotheses follow from Definition~\ref{def:measurement} alone. Those hypotheses are not proved; whether they hold for a given physical coupling is a separate modelling question. Alternative routes to sector separation may exist. The interaction functional $\Fcal_{\mathrm{int}}$ is introduced phenomenologically; its derivation from first principles or from the 600-cell geometry is deferred. The toy model below is illustrative, and will in fact illustrate a case in which the naive ``stronger coupling gives better separation'' intuition fails.

%==================================================================
\section{Toy Model: Coupled Double-Well Formation}\label{sec:toy}
%==================================================================

This model is illustrative. It is not derived from the full 600-cell construction, but demonstrates the landscape-deformation mechanism in the simplest possible setting.

\begin{example}[System--apparatus coupling]\label{ex:toy}
\textbf{System.} One-dimensional, with closure functional $\Fcal_S(\phi) = \alpha(\phi^2 - 1)^2$, having minima at $\phi = \pm 1$.

\textbf{Apparatus.} One-dimensional, with closure functional $\Fcal_A(\chi) = \tfrac{1}{2}\kappa \chi^2$, having a single minimum at $\chi = 0$ (the ``ready state'').

\textbf{Interaction.} $\Fcal_{\mathrm{int}}(\phi, \chi) = \tfrac{1}{2}\lambda(\chi - g\phi)^2$, favouring apparatus state $\chi = g\phi$ (the apparatus tracks the system).

\textbf{Total functional.}
\begin{equation}
    \Fcal_{\mathrm{total}}(\phi, \chi) = \alpha(\phi^2 - 1)^2 + \tfrac{1}{2}\kappa \chi^2 + \tfrac{1}{2}\lambda(\chi - g\phi)^2
\end{equation}

\textbf{Pre-measurement} ($\lambda = 0$). The system has two minima at $\phi = \pm 1$; the apparatus sits at $\chi = 0$. The total landscape has two minima at $(\pm 1, 0)$, with uncoupled-system barrier $\Delta\Fcal_0 = \alpha$ at $\phi = 0$. The apparatus configuration does not distinguish the two system minima.

\textbf{Post-measurement} ($\lambda > 0$). Minimising over $\chi$ at fixed $\phi$: $\chi^* = \lambda g\phi / (\kappa + \lambda)$. Substituting back gives the effective system functional
\begin{equation}
    \Fcal_{\mathrm{eff}}(\phi) \;=\; \alpha(\phi^2 - 1)^2 + c\,\phi^2, \qquad c \;=\; \frac{\kappa\lambda g^2}{2(\kappa + \lambda)}.
\end{equation}
Stationarity gives $\phi = 0$ or $\phi^2 = 1 - c/(2\alpha)$. The bistable structure \emph{survives only when $c < 2\alpha$}; for $c \geq 2\alpha$ the outer minima merge into a single minimum at $\phi = 0$ and the two-outcome interpretation is lost. Assuming $c < 2\alpha$, the correlated minima sit at
\begin{equation}
    (\phi, \chi)^{(1)*} \;\approx\; \bigl(+\sqrt{1 - c/(2\alpha)},\; +g\lambda\sqrt{1 - c/(2\alpha)}/(\kappa+\lambda)\bigr),
\end{equation}
and symmetrically at the negative root for $(\phi,\chi)^{(2)*}$. The apparatus coordinate $\chi^*(\phi)$ now distinguishes the two minima: $\chi > 0$ correlates with the $\phi > 0$ basin and $\chi < 0$ with the $\phi < 0$ basin.

\textbf{Barrier.} The saddle between the two minima sits at $(\phi, \chi) = (0, 0)$; direct computation gives
\begin{equation}
    \Delta\Fcal \;=\; \Fcal_{\mathrm{eff}}(0) - \Fcal_{\mathrm{eff}}(\phi^{(1)*}) \;=\; \alpha\left(1 - \frac{c}{2\alpha}\right)^2.
\end{equation}
This barrier \emph{decreases} with coupling in the bistable regime ($c < 2\alpha$): the uncoupled value $\alpha$ is the maximum, and the barrier vanishes at $c = 2\alpha$ where bistability is lost. Increasing the coupling strength $\lambda$ (or the tracking gain $g$) therefore \emph{weakens} the inter-basin barrier in this model, eventually destroying the two-outcome structure entirely.

\textbf{Sector separation (diagnostic, not conclusion).} For $\Delta\Fcal \gg \sigma^2$ the Kramers rate is $\Gamma \sim e^{-2\Delta\Fcal/\sigma^2}$; this regime is available only if $\alpha \gg \sigma^2$, and then only in a sufficiently small-$c$ window. Along the $\lambda$-varying family with fixed $\kappa$ and $g$, this corresponds to small $\lambda$; other parameter scalings can achieve small $c$ differently. The apparatus coordinate then records the system outcome by correlation. The lesson of the toy model, along the $\lambda$-varying family, is the reverse of a naive ``stronger coupling $\Rightarrow$ better separation'' intuition: the apparatus/system bistable architecture here is fragile under increasing $\lambda$ at fixed $\kappa, g$. A generic amplification mechanism that strengthens rather than weakens separation with coupling requires a different model (e.g.\ apparatus nonlinearities, pointer-basin intrinsic structure, or cascaded amplification), none of which is analysed here.

\textbf{Conclusion (for this toy model).} The system already possesses two minima, and the measurement interaction correlates them with the apparatus coordinate via $\chi^*(\phi) = \lambda g\phi/(\kappa + \lambda)$. The barrier formula $\Delta\Fcal = \alpha(1 - c/(2\alpha))^2$ shows that in the bistable regime the barrier \emph{decreases} with $c$ and vanishes at $c = 2\alpha$; and the pointer displacement at the matched minima is $\chi^*(\phi^*) = \lambda g \sqrt{1 - c/(2\alpha)}/(\kappa+\lambda)$. For \emph{fixed $\kappa$ and $g$}, varying $\lambda$ drives both $c$ and $\chi^*$: making $\lambda$ small enough to preserve a large barrier also keeps the pointer displacement small. Other parameter scalings (e.g.\ $\kappa \to 0$ at fixed $\lambda$) behave differently and can keep $\chi^*$ non-negligible while still having small $c$. The toy model therefore illustrates that, at least along the $\lambda$-varying family with other parameters fixed, barrier strength and readout amplitude trade off. It does not establish an unavoidable tension across all parameter regimes, and a robust measurement model remains to be constructed.
\end{example}

%==================================================================
\section{Connection to Paper~XXX}\label{sec:connection}
%==================================================================

Paper~XXX~\citep{SmartXXX} derived Born-like weighting $P_i = |c_i|^2$ under \emph{five} assumptions (equilibrium, well-separated sectors, sector-supported decomposition, a sampling postulate, and state-completeness). The present paper addresses, conditionally, one possible route toward the metastability part of assumption~2 of Paper~XXX; it does not establish that assumption as stated:

\begin{enumerate}
    \item \textbf{Well-separated sectors} (Assumption 2 of Paper~XXX). This paper offers a conditional mechanism template, via Remark~\ref{rem:basins}, by which Paper~XXX's sector-separation assumption could hold in a measurement context if the additional hypotheses (i)--(iv) of that remark are satisfied. No theorem that $\Delta\Fcal_{ij} \gg \sigma^2$ is proved here; the Kramers small-noise asymptotic~\eqref{eq:kramers} is imported from Paper~XIV.

    \item \textbf{Sector-supported decomposition} (Assumption 3 of Paper~XXX). Well-separated basins do \emph{not} by themselves produce a sector-supported basis; constructing localised modes $\phi_i$ with the support properties Paper~XXX uses requires a further spectral/semiclassical argument (Section~\ref{sec:suppression}), which is not carried out here.

    \item \textbf{Equilibrium} (Assumption 1 of Paper~XXX). Not addressed here. Paper~XIV proves a stationary distribution but does not prove generic convergence from arbitrary initial data; the equilibrium hypothesis remains conditional.

    \item \textbf{Sampling postulate and state-completeness} (Assumptions 4 and 5 of Paper~XXX). Not addressed here.
\end{enumerate}

\begin{quote}
This paper supplies a candidate sufficient-condition template (conditional on the hypotheses of Remark~\ref{rem:basins}) bearing on the metastability component of Paper~XXX's assumption~2 in a measurement context. The negligible-overlap component of assumption~2 and the sector-supported decomposition of assumption~3 remain entirely untouched here; four of Paper~XXX's five Route~C assumptions, including the overlap and decomposition components, remain external.
\end{quote}

%==================================================================
\section{Comparison to Standard Decoherence}\label{sec:comparison}
%==================================================================

\begin{table}[ht]
\centering
\caption{Structural comparison: closure sector separation vs.\ standard decoherence.}
\label{tab:comparison}
\begin{tabular}{@{}lll@{}}
\toprule
Feature & Standard decoherence & VFD sector separation \\
\midrule
Mechanism & Environmental entanglement & Proposed landscape barrier formation (not derived here) \\
Interference loss & Off-diagonal density matrix decay & Cross-term suppression from barrier expected heuristically (not proved here) \\
Pointer states & Einselection (environment-preferred) & Proposed basin minima of $\Fcal_{\mathrm{total}}$, if such minima form \\
Outcome stability & Decoherence timescale & Kramers barrier-crossing rate, conditional on Section~\ref{sec:suppression}'s three import hypotheses \\
Environment & Required (traced out) & Not required in this mechanism class \\
Formalism & Density matrices, partial trace & Closure functional, gradient flow \\
\bottomrule
\end{tabular}
\end{table}

\begin{remark}[Structural analogy, not equivalence]
The two frameworks are structurally analogous in aim; on the VFD side, sector separation is conditional on Remark~\ref{rem:basins}'s hypotheses and interference suppression remains heuristic (Section~\ref{sec:suppression}). The intended structural difference is that standard decoherence requires tracing over environmental degrees of freedom, while closure sector separation would, under the stated conditional template, arise from the geometry of the total closure landscape without invoking an external environment. Whether the two mechanisms are related by a deeper correspondence (e.g.\ whether environmental coupling can be recast as landscape deformation) is an open question.
\end{remark}

%==================================================================
\section{Discussion and Limitations}\label{sec:discussion}
%==================================================================

\subsection{What is established}

\begin{itemize}
    \item Measurement-type interactions are defined in terms of closure-landscape structure (Definition~\ref{def:measurement}).
    \item A heuristic template (Remark~\ref{rem:basins}) identifies four additional hypotheses --- (i) exhaustive minima, (ii) mismatch-crossing paths, (iii) a pairwise positive gap $\beta_{ij}>0$ over the mismatch set, and (iv) neighbourhood-level continuity/gap beyond Definition~\ref{def:measurement} --- under which the barrier lower bound~\eqref{eq:barrier} follows. Those hypotheses already assume most of the separation topology and are not derived from Definition~\ref{def:measurement}; the contribution is an accounting of what is assumed versus what follows, not a production mechanism for separated basins. Under the additional three Paper~XIV import hypotheses of Section~\ref{sec:suppression} (quasi-static landscape, pairwise saddle/nondegeneracy, two-basin reduction), the Kramers small-noise asymptotic then applies.
    \item If, in addition, the three Paper~XIV import hypotheses of Section~\ref{sec:suppression} hold (quasi-static landscape, pairwise saddle/nondegeneracy, two-basin reduction) \emph{and} $\Delta\Fcal_{i\to j} \gg \sigma^2$, the imported small-noise Kramers asymptotic (equation~\ref{eq:kramers}) would predict exponentially suppressed transitions.
    \item A toy model (Example~\ref{ex:toy}) shows that in a specific bistable system/apparatus setup the post-coupling barrier can be computed in closed form; within the bistable regime the barrier \emph{decreases} with the effective parameter $c$ and the bistable structure is destroyed at $c = 2\alpha$. If $\alpha \gg \sigma^2$, well-separated sectors are available only in a sufficiently small-$c$ window; along the fixed-$\kappa,g$, varying-$\lambda$ family this corresponds to small $\lambda$.
    \item A heuristic timescale hierarchy ansatz is stated (equation~\ref{eq:hierarchy}) as a working hypothesis for macroscopic apparatus.
\end{itemize}

\subsection{What is not established}

\begin{itemize}
    \item No proof that all physically relevant measurement interactions are of measurement type (Definition~\ref{def:measurement}).
    \item No proof of Remark~\ref{rem:basins}'s hypotheses: (i) exactly $N$ minima, (ii) every inter-minimum path crosses a mismatch configuration, (iii) the pairwise positive gap $\beta_{ij}>0$ of $\Fcal_{\mathrm{total}}$ over the mismatch set relative to the matched-minima baseline, and (iv) neighbourhood-level continuity/gap extending Definition~\ref{def:measurement}. All four are load-bearing for the barrier bound and none is implied by Definition~\ref{def:measurement} alone.
    \item No spectral/semiclassical construction of sector-localised modes $\phi_i$ with the support properties Paper~XXX uses; consequently, Paper~XXX's assumption~3 (sector-supported decomposition) is \emph{not} discharged.
    \item No proof of the overlap estimate $|\langle\phi_k|\phi_l\rangle| \lesssim e^{-c\Delta\Fcal/\sigma^2}$; this is only a schematic expectation.
    \item No derivation of $\Fcal_{\mathrm{int}}$ from first principles or from the 600-cell geometry.
    \item No proof of the timescale hierarchy~\eqref{eq:hierarchy} for general measurement contexts.
    \item No treatment of continuous-spectrum measurements or measurements with degenerate outcomes.
    \item No full comparison to standard decoherence theory (the comparison in Table~\ref{tab:comparison} is structural, not quantitative).
    \item No proof that the barrier mechanism is the \textit{only} route to sector separation; alternative mechanisms may exist.
    \item No generic amplification mechanism that strengthens rather than weakens separation with coupling; the toy model exhibits the opposite trend in the bistable regime.
    \item The remaining Paper~XXX Route~C assumptions --- equilibrium, sampling, and state-completeness --- are not addressed.
\end{itemize}

%==================================================================
\section{Conclusion}\label{sec:conclusion}
%==================================================================

Measurement, in the VFD framework proposed here, is recast, conditionally, as a barrier-based repackaging of sector separation rather than external state reduction. When a system couples to a measurement apparatus, the interaction term reshapes the closure functional; under the four additional hypotheses of Remark~\ref{rem:basins} (not proved here), the deformed landscape \emph{would} host $N$ apparatus-correlated basins (created by the coupling from a monostable pre-measurement landscape, or merely correlated with apparatus state if pre-measurement bistability was already present) whose directional barriers $\Delta\Fcal_{i\to j}$ satisfy the pairwise lower bound~\eqref{eq:barrier}. When these directional barriers satisfy $\Delta\Fcal_{i\to j} \gg \sigma^2$, the quasi-static post-interaction regime applies, the pairwise saddle/nondegeneracy hypotheses hold, and a justified two-basin reduction to the setting of Paper~XIV is available, the imported small-noise Kramers asymptotic would give exponentially suppressed transition rates.

This paper \emph{does not} discharge Paper~XXX's Route~C assumptions. It proposes a candidate mechanism class (conditional on Remark~\ref{rem:basins}'s hypotheses) by which Paper~XXX's sector-separation assumption could be licensed in a suitable measurement context, sketches the cross-term overlap estimate as a schematic expectation without proof, and works a toy model that \emph{fails} to demonstrate coupling-driven barrier amplification in its bistable regime. The other four of Paper~XXX's five Route~C assumptions --- equilibrium, sector-supported decomposition, sampling, and state-completeness --- are not addressed here.

\begin{enumerate}
    \item \textbf{This paper}: measurement interaction $\to$ \emph{candidate} landscape deformation $\to$ (under further hypotheses) sector separation.
    \item \textbf{Paper~XXX}: given well-separated sectors plus sampling plus completeness plus equilibrium plus sector-supported decomposition, the stationary sector measure yields Born weighting.
\end{enumerate}

The chain is conditional at every step. A complete, assumption-free measurement theory remains an open target. What is identified is a candidate mechanism class for one of Paper~XXX's assumptions, framed so that the quantitative work needed to convert it into a derivation can be undertaken in future work. Finding an amplification mechanism that strengthens rather than weakens separation with coupling --- needed if the mechanism class is to describe physical macroscopic measurements --- is the most immediate open problem.

%==================================================================
\bibliographystyle{plainnat}
\bibliography{references}

\end{document}
